
\documentclass{article}
\usepackage{colortbl}
\usepackage{makecell}
\usepackage{multirow}
\usepackage{supertabular}

\begin{document}

\newcounter{utterance}

\centering \large Interaction Transcript for game `clean\_up', experiment `2\_hard\_5obj\_lt', episode 3 with glm{-}5.2.
\vspace{24pt}

{ \footnotesize  \setcounter{utterance}{1}
\setlength{\tabcolsep}{0pt}
\begin{supertabular}{c@{$\;$}|p{.15\linewidth}@{}p{.15\linewidth}p{.15\linewidth}p{.15\linewidth}p{.15\linewidth}p{.15\linewidth}}
    \# & $\;$A & \multicolumn{4}{c}{Game Master} & $\;\:$B\\
    \hline

    \theutterance \stepcounter{utterance}  
    & & \multicolumn{4}{p{0.6\linewidth}}{
        \cellcolor[rgb]{0.9,0.9,0.9}{
            \makecell[{{p{\linewidth}}}]{
                \texttt{\tiny{[P1$\langle$GM]}}
                \texttt{Aš esu jūsų žaidimo vedėjas, o jūs žaidžiate bendrą žaidimą, kurio žaidimo lenta yra tokia lentelė:} \\
\\ 
\texttt{```} \\
\texttt{ 1234567} \\
\texttt{╔═╤═╤═══╗} \\
\texttt{║C│◌│◌TI║ 1} \\
\texttt{║◌├─┤W◌◌║ 2} \\
\texttt{╟─┤◌├─┬─╢ 3} \\
\texttt{║◌│◌│◌│◌║ 4} \\
\texttt{╟─┼─┘◌│◌║ 5} \\
\texttt{║◌│◌H┌┴─╢ 6} \\
\texttt{║◌│◌◌│◌◌║ 7} \\
\texttt{╚═╧══╧══╝} \\
\\ 
\texttt{```} \\
\\ 
\texttt{* Viršutiniame krašte x koordinatės didėja į dešinę, o dešiniajame krašte y koordinatės didėja žemyn.} \\
\texttt{* Tinklelyje atsitiktine tvarka išdėstyti šie objektai: 'W', 'I', 'T', 'C', 'H'.} \\
\\ 
\\ 
\texttt{Kitas žaidėjas mato kitą žaidimo lentos variantą, kuriame objektai išdėstyti skirtingose atsitiktinėse vietose. Jūs negalite matyti kito žaidėjo lentos, o jis negali matyti jūsų lentos.} \\
\\ 
\\ 
\texttt{**Tikslas:**} \\
\\ 
\texttt{Abu žaidėjai turi perkelti objektus savo fone taip, kad vienodi objektai atsidurtų tose pačiose koordinatėse. Turite bendrauti su kitu žaidėju ir susitarti dėl bendros tikslinės konfigūracijos.} \\
\\ 
\\ 
\texttt{**Taisyklės:**} \\
\\ 
\texttt{* Kiekvieno ėjimo metu galite išsiųsti lygiai vieną iš šių dviejų komandų:} \\
\texttt{ 1. `SAKYTI: <ŽINUTĖ>`: norėdami išsiųsti žinutę (viską iki kitos eilutės pertraukos) kitam žaidėjui. Aš ją persiųsiu jūsų partneriui.} \\
\texttt{2. `JUDĖTI: <OBJEKTAS>, (<X>, <Y>)`: perkelti objektą į naują poziciją, kur `<X>` yra stulpelis, o `<Y>` yra eilutė. Aš jums pranešiu, ar jūsų perkėlimas buvo sėkmingas, ar ne.} \\
\texttt{* Jei nesilaikysite formato arba siųsite kelias komandas iš karto, turėsiu jus nubausti.} \\
\texttt{* Jei abu žaidėjai sukaups daugiau nei 12 baudų, abu pralaimi žaidimą.} \\
\texttt{* Labai svarbu, kad bendrautumėte su kitu žaidėju dėl savo tikslinės būsenos! Vienintelis būdas, kuriuo galite perduoti savo strategiją kitam žaidėjui, yra komanda `SAKYTI: <ŽINUTĖ>`!} \\
\\ 
\\ 
\texttt{**Objektų perkėlimas**} \\
\\ 
\texttt{* Objektus galima perkelti tik į langelius, esančius tinklelio ribose. Tikslinis langelis turi būti tuščias, t. y. jame turi būti tik simbolis '◌'.} \\
\texttt{* Jei bandysite perkelti objektą į vietą, kuri nėra tuščia, arba bandysite jį perkelti už tinklelio ribų, turėsiu jus nubausti. Gausite dar vieną bandymą.} \\
\texttt{* Prieš atlikdami perkėlimą, dukart patikrinkite, ar tikslinė vieta yra tuščia ir ar joje nėra jokios raidės, rėmelio ar linijos!} \\
\\ 
\\ 
\texttt{**Žaidimo pabaiga**} \\
\\ 
\texttt{Jei manote, kad pasiekėte tikslą {-} sulygiavote visus objektus, galite paprašyti kito žaidėjo baigti žaidimą siųsdami `SAKYTI: baigta?`. Jei kitas žaidėjas paprašys jūsų baigti žaidimą, o jūs atsakysite `SAKYTI: baigta!`, žaidimas baigsis.} \\
\\ 
\texttt{Abu žaidėjai laimi, jei žaidimas baigiasi per 20 raundus, kur vienas raundas apibrėžiamas kaip du žaidėjai, kiekvienas išsiuntęs po vieną galiojančią komandą.} \\
\\ 
\texttt{**Taškų skaičiavimas:**} \\
\\ 
\texttt{Kuo arčiau vienodų objektų abiejose žaidimo lentose, tuo daugiau taškų gausite. Baudos mažina jūsų taškus. Ar pavyks pagerinti rekordą?} \\
\\ 
\\ 
\texttt{Išsiųskite žinutę kitam žaidėjui, kad pradėtumėte žaidimą!} \\
            }
        }
    }
    & \\ \\

    \theutterance \stepcounter{utterance}  
    & \multicolumn{4}{p{0.6\linewidth}}{
        \cellcolor[rgb]{0.9,0.9,0.9}{
            \makecell[{{p{\linewidth}}}]{
                \texttt{\tiny{[P1$\rangle$GM]}}
                \texttt{SAKYTI: Sveiki! Pasiūlymas dėl tikslinių koordinacžių (x, y): C=(3,4), H=(5,4), I=(5,1), T=(4,1), W=(2,3). Ar tinka?} \\
            }
        }
    }
    & & \\ \\

    \theutterance \stepcounter{utterance}  
    & & & \multicolumn{2}{p{0.3\linewidth}}{
        \cellcolor[rgb]{0.9,0.9,0.9}{
            \makecell[{{p{\linewidth}}}]{
                \texttt{\tiny{[GM$|$GM]}}
                \texttt{SAKYTI: Sveiki! Pasiūlymas dėl tikslinių koordinacžių (x, y): C=(3,4), H=(5,4), I=(5,1), T=(4,1), W=(2,3). Ar tinka?} \\
            }
        }
    }
    & & \\ \\

    \theutterance \stepcounter{utterance}  
    & & \multicolumn{4}{p{0.6\linewidth}}{
        \cellcolor[rgb]{0.9,0.9,0.9}{
            \makecell[{{p{\linewidth}}}]{
                \texttt{\tiny{[GM$\rangle$P2]}}
                \texttt{Aš esu jūsų žaidimo vedėjas, o jūs žaidžiate bendrą žaidimą, kurio žaidimo lenta yra tokia lentelė:} \\
\\ 
\texttt{```} \\
\texttt{ 1234567} \\
\texttt{╔═╤═╤═══╗} \\
\texttt{║H│◌│◌◌◌║ 1} \\
\texttt{║◌├─┤◌◌◌║ 2} \\
\texttt{╟─┤◌├─┬─╢ 3} \\
\texttt{║◌│◌│C│◌║ 4} \\
\texttt{╟─┼─┘◌│◌║ 5} \\
\texttt{║◌│◌◌┌┴─╢ 6} \\
\texttt{║T│WI│◌◌║ 7} \\
\texttt{╚═╧══╧══╝} \\
\\ 
\texttt{```} \\
\\ 
\texttt{* Viršutiniame krašte x koordinatės didėja į dešinę, o dešiniajame krašte y koordinatės didėja žemyn.} \\
\texttt{* Tinklelyje atsitiktine tvarka išdėstyti šie objektai: 'W', 'I', 'T', 'C', 'H'.} \\
\\ 
\\ 
\texttt{Kitas žaidėjas mato kitą žaidimo lentos variantą, kuriame objektai išdėstyti skirtingose atsitiktinėse vietose. Jūs negalite matyti kito žaidėjo lentos, o jis negali matyti jūsų lentos.} \\
\\ 
\\ 
\texttt{**Tikslas:**} \\
\\ 
\texttt{Abu žaidėjai turi perkelti objektus savo fone taip, kad vienodi objektai atsidurtų tose pačiose koordinatėse. Turite bendrauti su kitu žaidėju ir susitarti dėl bendros tikslinės konfigūracijos.} \\
\\ 
\\ 
\texttt{**Taisyklės:**} \\
\\ 
\texttt{* Kiekvieno ėjimo metu galite išsiųsti lygiai vieną iš šių dviejų komandų:} \\
\texttt{ 1. `SAKYTI: <ŽINUTĖ>`: norėdami išsiųsti žinutę (viską iki kitos eilutės pertraukos) kitam žaidėjui. Aš ją persiųsiu jūsų partneriui.} \\
\texttt{2. `JUDĖTI: <OBJEKTAS>, (<X>, <Y>)`: perkelti objektą į naują poziciją, kur `<X>` yra stulpelis, o `<Y>` yra eilutė. Aš jums pranešiu, ar jūsų perkėlimas buvo sėkmingas, ar ne.} \\
\texttt{* Jei nesilaikysite formato arba siųsite kelias komandas iš karto, turėsiu jus nubausti.} \\
\texttt{* Jei abu žaidėjai sukaups daugiau nei 12 baudų, abu pralaimi žaidimą.} \\
\texttt{* Labai svarbu, kad bendrautumėte su kitu žaidėju dėl savo tikslinės būsenos! Vienintelis būdas, kuriuo galite perduoti savo strategiją kitam žaidėjui, yra komanda `SAKYTI: <ŽINUTĖ>`!} \\
\\ 
\\ 
\texttt{**Objektų perkėlimas**} \\
\\ 
\texttt{* Objektus galima perkelti tik į langelius, esančius tinklelio ribose. Tikslinis langelis turi būti tuščias, t. y. jame turi būti tik simbolis '◌'.} \\
\texttt{* Jei bandysite perkelti objektą į vietą, kuri nėra tuščia, arba bandysite jį perkelti už tinklelio ribų, turėsiu jus nubausti. Gausite dar vieną bandymą.} \\
\texttt{* Prieš atlikdami perkėlimą, dukart patikrinkite, ar tikslinė vieta yra tuščia ir ar joje nėra jokios raidės, rėmelio ar linijos!} \\
\\ 
\\ 
\texttt{**Žaidimo pabaiga**} \\
\\ 
\texttt{Jei manote, kad pasiekėte tikslą {-} sulygiavote visus objektus, galite paprašyti kito žaidėjo baigti žaidimą siųsdami `SAKYTI: baigta?`. Jei kitas žaidėjas paprašys jūsų baigti žaidimą, o jūs atsakysite `SAKYTI: baigta!`, žaidimas baigsis.} \\
\\ 
\texttt{Abu žaidėjai laimi, jei žaidimas baigiasi per 20 raundus, kur vienas raundas apibrėžiamas kaip du žaidėjai, kiekvienas išsiuntęs po vieną galiojančią komandą.} \\
\\ 
\texttt{**Taškų skaičiavimas:**} \\
\\ 
\texttt{Kuo arčiau vienodų objektų abiejose žaidimo lentose, tuo daugiau taškų gausite. Baudos mažina jūsų taškus. Ar pavyks pagerinti rekordą?} \\
\\ 
\texttt{Kitas žaidėjas pradėjo žaidimą išsiųsdamas šią žinutę:} \\
\\ 
\texttt{"Sveiki! Pasiūlymas dėl tikslinių koordinacžių (x, y): C=(3,4), H=(5,4), I=(5,1), T=(4,1), W=(2,3). Ar tinka?"} \\
\\ 
\texttt{Kokia yra jūsų pirmoji komanda?} \\
            }
        }
    }
    & \\ \\

    \theutterance \stepcounter{utterance}  
    & & & \multicolumn{4}{p{0.6\linewidth}}{
        \cellcolor[rgb]{0.9,0.9,0.9}{
            \makecell[{{p{\linewidth}}}]{
                \texttt{\tiny{[GM$\langle$P2]}}
                \texttt{SAKYTI: Sveiki! Deja, negaliu sutikti su tomis koordinatėmis, nes kai kurios jos mano lentoje yra užimtos (pvz., C=(3,4) jau yra H=(5,4) vietoje tiesiogiai nepasiekiama be kitų perkėlimų, o (4,1) užimta). Siūlau tokias laisvas koordinates: H=(1,2), I=(1,5), T=(7,5), C=(3,4), W=(5,5). Ar tinka?} \\
            }
        }
    }
     \\ \\

    \theutterance \stepcounter{utterance}  
    & & & \multicolumn{2}{p{0.3\linewidth}}{
        \cellcolor[rgb]{0.9,0.9,0.9}{
            \makecell[{{p{\linewidth}}}]{
                \texttt{\tiny{[GM$|$GM]}}
                \texttt{SAKYTI: Sveiki! Deja, negaliu sutikti su tomis koordinatėmis, nes kai kurios jos mano lentoje yra užimtos (pvz., C=(3,4) jau yra H=(5,4) vietoje tiesiogiai nepasiekiama be kitų perkėlimų, o (4,1) užimta). Siūlau tokias laisvas koordinates: H=(1,2), I=(1,5), T=(7,5), C=(3,4), W=(5,5). Ar tinka?} \\
            }
        }
    }
    & & \\ \\

    \theutterance \stepcounter{utterance}  
    & & \multicolumn{4}{p{0.6\linewidth}}{
        \cellcolor[rgb]{0.9,0.9,0.9}{
            \makecell[{{p{\linewidth}}}]{
                \texttt{\tiny{[P1$\langle$GM]}}
                \texttt{Jūsų žinutė buvo perduota kitam žaidėjui.} \\
\\ 
\\ 
\texttt{Šiuo metu žaidžiate 1 raundą iš maksimaliai 20.} \\
\texttt{Jūs bendrai sukaupėte 0 iš 12 baudų.} \\
\texttt{Kitas žaidėjas atsiuntė šią žinutę:} \\
\\ 
\\ 
\texttt{"Sveiki! Deja, negaliu sutikti su tomis koordinatėmis, nes kai kurios jos mano lentoje yra užimtos (pvz., C=(3,4) jau yra H=(5,4) vietoje tiesiogiai nepasiekiama be kitų perkėlimų, o (4,1) užimta). Siūlau tokias laisvas koordinates: H=(1,2), I=(1,5), T=(7,5), C=(3,4), W=(5,5). Ar tinka?"} \\
\\ 
\\ 
\texttt{Kokia jūsų kita komanda?} \\
            }
        }
    }
    & \\ \\

    \theutterance \stepcounter{utterance}  
    & \multicolumn{4}{p{0.6\linewidth}}{
        \cellcolor[rgb]{0.9,0.9,0.9}{
            \makecell[{{p{\linewidth}}}]{
                \texttt{\tiny{[P1$\rangle$GM]}}
                \texttt{SAKYTI: Tinka! Susitariame: H=(1,2), I=(1,5), T=(7,5), C=(3,4), W=(5,5). Pradedu perkėlimą.} \\
            }
        }
    }
    & & \\ \\

    \theutterance \stepcounter{utterance}  
    & & & \multicolumn{2}{p{0.3\linewidth}}{
        \cellcolor[rgb]{0.9,0.9,0.9}{
            \makecell[{{p{\linewidth}}}]{
                \texttt{\tiny{[GM$|$GM]}}
                \texttt{SAKYTI: Tinka! Susitariame: H=(1,2), I=(1,5), T=(7,5), C=(3,4), W=(5,5). Pradedu perkėlimą.} \\
            }
        }
    }
    & & \\ \\

    \theutterance \stepcounter{utterance}  
    & & \multicolumn{4}{p{0.6\linewidth}}{
        \cellcolor[rgb]{0.9,0.9,0.9}{
            \makecell[{{p{\linewidth}}}]{
                \texttt{\tiny{[GM$\rangle$P2]}}
                \texttt{Jūsų žinutė buvo perduota kitam žaidėjui.} \\
\\ 
\\ 
\texttt{Šiuo metu žaidžiate 2 raundą iš maksimaliai 20.} \\
\texttt{Jūs bendrai sukaupėte 0 iš 12 baudų.} \\
\texttt{Kitas žaidėjas atsiuntė šią žinutę:} \\
\\ 
\\ 
\texttt{"Tinka! Susitariame: H=(1,2), I=(1,5), T=(7,5), C=(3,4), W=(5,5). Pradedu perkėlimą."} \\
\\ 
\\ 
\texttt{Kokia jūsų kita komanda?} \\
            }
        }
    }
    & \\ \\

    \theutterance \stepcounter{utterance}  
    & & & \multicolumn{4}{p{0.6\linewidth}}{
        \cellcolor[rgb]{0.9,0.9,0.9}{
            \makecell[{{p{\linewidth}}}]{
                \texttt{\tiny{[GM$\langle$P2]}}
                \texttt{JUDĖTI: H, (1, 2)} \\
            }
        }
    }
     \\ \\

    \theutterance \stepcounter{utterance}  
    & & & \multicolumn{2}{p{0.3\linewidth}}{
        \cellcolor[rgb]{0.9,0.9,0.9}{
            \makecell[{{p{\linewidth}}}]{
                \texttt{\tiny{[GM$|$GM]}}
                \texttt{JUDĖTI: H, (1, 2)} \\
            }
        }
    }
    & & \\ \\

    \theutterance \stepcounter{utterance}  
    & & & \multicolumn{2}{p{0.3\linewidth}}{
        \cellcolor[rgb]{0.9,0.9,0.9}{
            \makecell[{{p{\linewidth}}}]{
                \texttt{\tiny{[GM$|$GM]}}
                \texttt{H moved to (1, 2)} \\
            }
        }
    }
    & & \\ \\

    \theutterance \stepcounter{utterance}  
    & & \multicolumn{4}{p{0.6\linewidth}}{
        \cellcolor[rgb]{0.9,0.9,0.9}{
            \makecell[{{p{\linewidth}}}]{
                \texttt{\tiny{[P1$\langle$GM]}}
                \texttt{Jūsų žinutė buvo perduota kitam žaidėjui.} \\
\\ 
\\ 
\texttt{Šiuo metu žaidžiate 2 raundą iš maksimaliai 20.} \\
\texttt{Jūs bendrai sukaupėte 0 iš 12 baudų.} \\
\texttt{Kitas žaidėjas perkėlė objektą savo tinklelyje.} \\
\\ 
\\ 
\texttt{Kokia jūsų kita komanda?} \\
            }
        }
    }
    & \\ \\

    \theutterance \stepcounter{utterance}  
    & \multicolumn{4}{p{0.6\linewidth}}{
        \cellcolor[rgb]{0.9,0.9,0.9}{
            \makecell[{{p{\linewidth}}}]{
                \texttt{\tiny{[P1$\rangle$GM]}}
                \texttt{Pirma patikrinsiu savo lentą. Mano objektų pradinės pozicijos:} \\
\texttt{{-} C: (1,1)} \\
\texttt{{-} T: (5,1)} \\
\texttt{{-} I: (7,1)} \\
\texttt{{-} W: (3,2)} \\
\texttt{{-} H: (4,6)} \\
\\ 
\texttt{Tikslinės pozicijos: H=(1,2), I=(1,5), T=(7,5), C=(3,4), W=(5,5).} \\
\\ 
\texttt{Pradėsiu perkeldamas H į (1,2), nes (1,2) yra tuščias langelis.} \\
\\ 
\texttt{JUDĖTI: H, (1,2)} \\
            }
        }
    }
    & & \\ \\

    \theutterance \stepcounter{utterance}  
    & & & \multicolumn{2}{p{0.3\linewidth}}{
        \cellcolor[rgb]{0.9,0.9,0.9}{
            \makecell[{{p{\linewidth}}}]{
                \texttt{\tiny{[GM$|$GM]}}
                \texttt{Pirma patikrinsiu savo lentą. Mano objektų pradinės pozicijos:} \\
\texttt{{-} C: (1,1)} \\
\texttt{{-} T: (5,1)} \\
\texttt{{-} I: (7,1)} \\
\texttt{{-} W: (3,2)} \\
\texttt{{-} H: (4,6)} \\
\\ 
\texttt{Tikslinės pozicijos: H=(1,2), I=(1,5), T=(7,5), C=(3,4), W=(5,5).} \\
\\ 
\texttt{Pradėsiu perkeldamas H į (1,2), nes (1,2) yra tuščias langelis.} \\
\\ 
\texttt{JUDĖTI: H, (1,2)} \\
            }
        }
    }
    & & \\ \\

    \theutterance \stepcounter{utterance}  
    & & & \multicolumn{2}{p{0.3\linewidth}}{
        \cellcolor[rgb]{0.9,0.9,0.9}{
            \makecell[{{p{\linewidth}}}]{
                \texttt{\tiny{[GM$|$GM]}}
                \texttt{Invalid format: head is not empty: 'Pirma patikrinsiu savo lentą. Mano objektų pradinės pozicijos:} \\
\texttt{{-} C: (1,1)} \\
\texttt{{-} T: (5,1)} \\
\texttt{{-} I: (7,1)} \\
\texttt{{-} W: (3,2)} \\
\texttt{{-} H: (4,6)} \\
\\ 
\texttt{Tikslinės pozicijos: H=(1,2), I=(1,5), T=(7,5), C=(3,4), W=(5,5).} \\
\\ 
\texttt{Pradėsiu perkeldamas H į (1,2), nes (1,2) yra tuščias langelis.} \\
\\ 
\texttt{'} \\
            }
        }
    }
    & & \\ \\

    \theutterance \stepcounter{utterance}  
    & & \multicolumn{4}{p{0.6\linewidth}}{
        \cellcolor[rgb]{0.9,0.9,0.9}{
            \makecell[{{p{\linewidth}}}]{
                \texttt{\tiny{[P1$\langle$GM]}}
                \texttt{Nuobauda: Jūsų žinutėje neturi būti nieko prieš komandą!} \\
\\ 
\texttt{Įsitikinkite, kad jūsų atsakyme yra tik `SAKYTI: <ŽINUTĖ>` arba `JUDĖTI: <OBJEKTAS>, (<X>, <Y>)`, ir nieko daugiau!} \\
\\ 
\texttt{Jūs bendrai sukaupėte 1 iš 12 baudų.} \\
\texttt{Pabandykite dar kartą!} \\
            }
        }
    }
    & \\ \\

    \theutterance \stepcounter{utterance}  
    & \multicolumn{4}{p{0.6\linewidth}}{
        \cellcolor[rgb]{0.9,0.9,0.9}{
            \makecell[{{p{\linewidth}}}]{
                \texttt{\tiny{[P1$\rangle$GM]}}
                \texttt{JUDĖTI: H, (1,2)} \\
            }
        }
    }
    & & \\ \\

    \theutterance \stepcounter{utterance}  
    & & & \multicolumn{2}{p{0.3\linewidth}}{
        \cellcolor[rgb]{0.9,0.9,0.9}{
            \makecell[{{p{\linewidth}}}]{
                \texttt{\tiny{[GM$|$GM]}}
                \texttt{JUDĖTI: H, (1,2)} \\
            }
        }
    }
    & & \\ \\

    \theutterance \stepcounter{utterance}  
    & & & \multicolumn{2}{p{0.3\linewidth}}{
        \cellcolor[rgb]{0.9,0.9,0.9}{
            \makecell[{{p{\linewidth}}}]{
                \texttt{\tiny{[GM$|$GM]}}
                \texttt{H moved to (1, 2)} \\
            }
        }
    }
    & & \\ \\

    \theutterance \stepcounter{utterance}  
    & & \multicolumn{4}{p{0.6\linewidth}}{
        \cellcolor[rgb]{0.9,0.9,0.9}{
            \makecell[{{p{\linewidth}}}]{
                \texttt{\tiny{[GM$\rangle$P2]}}
                \texttt{Sėkmingai perkelta 'H' į (1,2). Jūsų atnaujintas tinklelis atrodo taip:} \\
\texttt{```} \\
\texttt{ 1234567} \\
\texttt{╔═╤═╤═══╗} \\
\texttt{║◌│◌│◌◌◌║ 1} \\
\texttt{║H├─┤◌◌◌║ 2} \\
\texttt{╟─┤◌├─┬─╢ 3} \\
\texttt{║◌│◌│C│◌║ 4} \\
\texttt{╟─┼─┘◌│◌║ 5} \\
\texttt{║◌│◌◌┌┴─╢ 6} \\
\texttt{║T│WI│◌◌║ 7} \\
\texttt{╚═╧══╧══╝} \\
\texttt{```} \\
\texttt{Šiuo metu žaidžiate 3 raundą iš maksimaliai 20.} \\
\texttt{Jūs bendrai sukaupėte 1 iš 12 baudų.} \\
\texttt{Kitas žaidėjas perkėlė objektą savo tinklelyje.} \\
\\ 
\\ 
\texttt{Kokia jūsų kita komanda?} \\
            }
        }
    }
    & \\ \\

    \theutterance \stepcounter{utterance}  
    & & & \multicolumn{4}{p{0.6\linewidth}}{
        \cellcolor[rgb]{0.9,0.9,0.9}{
            \makecell[{{p{\linewidth}}}]{
                \texttt{\tiny{[GM$\langle$P2]}}
                \texttt{SAKYTI: Puiku, aš perkėliau H į (1,2). Dabar kelsiu W iš (2,7) į (5,5). Kitas žingsnis bus I į (1,5). Kaip sekasi tau?} \\
            }
        }
    }
     \\ \\

    \theutterance \stepcounter{utterance}  
    & & & \multicolumn{2}{p{0.3\linewidth}}{
        \cellcolor[rgb]{0.9,0.9,0.9}{
            \makecell[{{p{\linewidth}}}]{
                \texttt{\tiny{[GM$|$GM]}}
                \texttt{SAKYTI: Puiku, aš perkėliau H į (1,2). Dabar kelsiu W iš (2,7) į (5,5). Kitas žingsnis bus I į (1,5). Kaip sekasi tau?} \\
            }
        }
    }
    & & \\ \\

    \theutterance \stepcounter{utterance}  
    & & \multicolumn{4}{p{0.6\linewidth}}{
        \cellcolor[rgb]{0.9,0.9,0.9}{
            \makecell[{{p{\linewidth}}}]{
                \texttt{\tiny{[P1$\langle$GM]}}
                \texttt{Sėkmingai perkelta 'H' į (1,2). Jūsų atnaujintas tinklelis atrodo taip:} \\
\texttt{```} \\
\texttt{ 1234567} \\
\texttt{╔═╤═╤═══╗} \\
\texttt{║C│◌│◌TI║ 1} \\
\texttt{║H├─┤W◌◌║ 2} \\
\texttt{╟─┤◌├─┬─╢ 3} \\
\texttt{║◌│◌│◌│◌║ 4} \\
\texttt{╟─┼─┘◌│◌║ 5} \\
\texttt{║◌│◌◌┌┴─╢ 6} \\
\texttt{║◌│◌◌│◌◌║ 7} \\
\texttt{╚═╧══╧══╝} \\
\texttt{```} \\
\texttt{Šiuo metu žaidžiate 3 raundą iš maksimaliai 20.} \\
\texttt{Jūs bendrai sukaupėte 1 iš 12 baudų.} \\
\texttt{Kitas žaidėjas atsiuntė šią žinutę:} \\
\\ 
\\ 
\texttt{"Puiku, aš perkėliau H į (1,2). Dabar kelsiu W iš (2,7) į (5,5). Kitas žingsnis bus I į (1,5). Kaip sekasi tau?"} \\
\\ 
\\ 
\texttt{Kokia jūsų kita komanda?} \\
            }
        }
    }
    & \\ \\

    \theutterance \stepcounter{utterance}  
    & \multicolumn{4}{p{0.6\linewidth}}{
        \cellcolor[rgb]{0.9,0.9,0.9}{
            \makecell[{{p{\linewidth}}}]{
                \texttt{\tiny{[P1$\rangle$GM]}}
                \texttt{JUDĖTI: W, (5,5)} \\
            }
        }
    }
    & & \\ \\

    \theutterance \stepcounter{utterance}  
    & & & \multicolumn{2}{p{0.3\linewidth}}{
        \cellcolor[rgb]{0.9,0.9,0.9}{
            \makecell[{{p{\linewidth}}}]{
                \texttt{\tiny{[GM$|$GM]}}
                \texttt{JUDĖTI: W, (5,5)} \\
            }
        }
    }
    & & \\ \\

    \theutterance \stepcounter{utterance}  
    & & & \multicolumn{2}{p{0.3\linewidth}}{
        \cellcolor[rgb]{0.9,0.9,0.9}{
            \makecell[{{p{\linewidth}}}]{
                \texttt{\tiny{[GM$|$GM]}}
                \texttt{W moved to (5, 5)} \\
            }
        }
    }
    & & \\ \\

    \theutterance \stepcounter{utterance}  
    & & \multicolumn{4}{p{0.6\linewidth}}{
        \cellcolor[rgb]{0.9,0.9,0.9}{
            \makecell[{{p{\linewidth}}}]{
                \texttt{\tiny{[GM$\rangle$P2]}}
                \texttt{Jūsų žinutė buvo perduota kitam žaidėjui.} \\
\\ 
\\ 
\texttt{Šiuo metu žaidžiate 4 raundą iš maksimaliai 20.} \\
\texttt{Jūs bendrai sukaupėte 1 iš 12 baudų.} \\
\texttt{Kitas žaidėjas perkėlė objektą savo tinklelyje.} \\
\\ 
\\ 
\texttt{Kokia jūsų kita komanda?} \\
            }
        }
    }
    & \\ \\

    \theutterance \stepcounter{utterance}  
    & & & \multicolumn{4}{p{0.6\linewidth}}{
        \cellcolor[rgb]{0.9,0.9,0.9}{
            \makecell[{{p{\linewidth}}}]{
                \texttt{\tiny{[GM$\langle$P2]}}
                \texttt{JUDĖTI: W, (5, 5)} \\
            }
        }
    }
     \\ \\

    \theutterance \stepcounter{utterance}  
    & & & \multicolumn{2}{p{0.3\linewidth}}{
        \cellcolor[rgb]{0.9,0.9,0.9}{
            \makecell[{{p{\linewidth}}}]{
                \texttt{\tiny{[GM$|$GM]}}
                \texttt{JUDĖTI: W, (5, 5)} \\
            }
        }
    }
    & & \\ \\

    \theutterance \stepcounter{utterance}  
    & & & \multicolumn{2}{p{0.3\linewidth}}{
        \cellcolor[rgb]{0.9,0.9,0.9}{
            \makecell[{{p{\linewidth}}}]{
                \texttt{\tiny{[GM$|$GM]}}
                \texttt{W moved to (5, 5)} \\
            }
        }
    }
    & & \\ \\

    \theutterance \stepcounter{utterance}  
    & & \multicolumn{4}{p{0.6\linewidth}}{
        \cellcolor[rgb]{0.9,0.9,0.9}{
            \makecell[{{p{\linewidth}}}]{
                \texttt{\tiny{[P1$\langle$GM]}}
                \texttt{Sėkmingai perkelta 'W' į (5,5). Jūsų atnaujintas tinklelis atrodo taip:} \\
\texttt{```} \\
\texttt{ 1234567} \\
\texttt{╔═╤═╤═══╗} \\
\texttt{║C│◌│◌TI║ 1} \\
\texttt{║H├─┤◌◌◌║ 2} \\
\texttt{╟─┤◌├─┬─╢ 3} \\
\texttt{║◌│◌│◌│◌║ 4} \\
\texttt{╟─┼─┘W│◌║ 5} \\
\texttt{║◌│◌◌┌┴─╢ 6} \\
\texttt{║◌│◌◌│◌◌║ 7} \\
\texttt{╚═╧══╧══╝} \\
\texttt{```} \\
\texttt{Šiuo metu žaidžiate 4 raundą iš maksimaliai 20.} \\
\texttt{Jūs bendrai sukaupėte 1 iš 12 baudų.} \\
\texttt{Kitas žaidėjas perkėlė objektą savo tinklelyje.} \\
\\ 
\\ 
\texttt{Kokia jūsų kita komanda?} \\
            }
        }
    }
    & \\ \\

    \theutterance \stepcounter{utterance}  
    & \multicolumn{4}{p{0.6\linewidth}}{
        \cellcolor[rgb]{0.9,0.9,0.9}{
            \makecell[{{p{\linewidth}}}]{
                \texttt{\tiny{[P1$\rangle$GM]}}
                \texttt{JUDĖTI: I, (1,5)} \\
            }
        }
    }
    & & \\ \\

    \theutterance \stepcounter{utterance}  
    & & & \multicolumn{2}{p{0.3\linewidth}}{
        \cellcolor[rgb]{0.9,0.9,0.9}{
            \makecell[{{p{\linewidth}}}]{
                \texttt{\tiny{[GM$|$GM]}}
                \texttt{JUDĖTI: I, (1,5)} \\
            }
        }
    }
    & & \\ \\

    \theutterance \stepcounter{utterance}  
    & & & \multicolumn{2}{p{0.3\linewidth}}{
        \cellcolor[rgb]{0.9,0.9,0.9}{
            \makecell[{{p{\linewidth}}}]{
                \texttt{\tiny{[GM$|$GM]}}
                \texttt{Nuobauda: (1,5) nėra tuščias, bet jame yra '─'.} \\
\\ 
\texttt{Jūs bendrai sukaupėte 2 iš 12 baudų.} \\
\texttt{Pabandykite dar kartą!} \\
            }
        }
    }
    & & \\ \\

    \theutterance \stepcounter{utterance}  
    & & \multicolumn{4}{p{0.6\linewidth}}{
        \cellcolor[rgb]{0.9,0.9,0.9}{
            \makecell[{{p{\linewidth}}}]{
                \texttt{\tiny{[P1$\langle$GM]}}
                \texttt{Nuobauda: (1,5) nėra tuščias, bet jame yra '─'.} \\
\\ 
\texttt{Jūs bendrai sukaupėte 2 iš 12 baudų.} \\
\texttt{Pabandykite dar kartą!} \\
            }
        }
    }
    & \\ \\

    \theutterance \stepcounter{utterance}  
    & \multicolumn{4}{p{0.6\linewidth}}{
        \cellcolor[rgb]{0.9,0.9,0.9}{
            \makecell[{{p{\linewidth}}}]{
                \texttt{\tiny{[P1$\rangle$GM]}}
                \texttt{SAKYTI: (1,5) turi rėmelio liniją. Ar (7,5) pas tave laisvas? Reikia rasti kitas laisvas koordinates I ir T. Mano laisvi langeliai: (1,4),(2,4),(3,4),(4,4),(7,4),(1,7),(2,7),(4,7),(5,7),(6,7),(7,7).} \\
            }
        }
    }
    & & \\ \\

    \theutterance \stepcounter{utterance}  
    & & & \multicolumn{2}{p{0.3\linewidth}}{
        \cellcolor[rgb]{0.9,0.9,0.9}{
            \makecell[{{p{\linewidth}}}]{
                \texttt{\tiny{[GM$|$GM]}}
                \texttt{SAKYTI: (1,5) turi rėmelio liniją. Ar (7,5) pas tave laisvas? Reikia rasti kitas laisvas koordinates I ir T. Mano laisvi langeliai: (1,4),(2,4),(3,4),(4,4),(7,4),(1,7),(2,7),(4,7),(5,7),(6,7),(7,7).} \\
            }
        }
    }
    & & \\ \\

    \theutterance \stepcounter{utterance}  
    & & \multicolumn{4}{p{0.6\linewidth}}{
        \cellcolor[rgb]{0.9,0.9,0.9}{
            \makecell[{{p{\linewidth}}}]{
                \texttt{\tiny{[GM$\rangle$P2]}}
                \texttt{Sėkmingai perkelta 'W' į (5,5). Jūsų atnaujintas tinklelis atrodo taip:} \\
\texttt{```} \\
\texttt{ 1234567} \\
\texttt{╔═╤═╤═══╗} \\
\texttt{║◌│◌│◌◌◌║ 1} \\
\texttt{║H├─┤◌◌◌║ 2} \\
\texttt{╟─┤◌├─┬─╢ 3} \\
\texttt{║◌│◌│C│◌║ 4} \\
\texttt{╟─┼─┘W│◌║ 5} \\
\texttt{║◌│◌◌┌┴─╢ 6} \\
\texttt{║T│◌I│◌◌║ 7} \\
\texttt{╚═╧══╧══╝} \\
\texttt{```} \\
\texttt{Šiuo metu žaidžiate 5 raundą iš maksimaliai 20.} \\
\texttt{Jūs bendrai sukaupėte 2 iš 12 baudų.} \\
\texttt{Kitas žaidėjas atsiuntė šią žinutę:} \\
\\ 
\\ 
\texttt{"(1,5) turi rėmelio liniją. Ar (7,5) pas tave laisvas? Reikia rasti kitas laisvas koordinates I ir T. Mano laisvi langeliai: (1,4),(2,4),(3,4),(4,4),(7,4),(1,7),(2,7),(4,7),(5,7),(6,7),(7,7)."} \\
\\ 
\\ 
\texttt{Kokia jūsų kita komanda?} \\
            }
        }
    }
    & \\ \\

    \theutterance \stepcounter{utterance}  
    & & & \multicolumn{4}{p{0.6\linewidth}}{
        \cellcolor[rgb]{0.9,0.9,0.9}{
            \makecell[{{p{\linewidth}}}]{
                \texttt{\tiny{[GM$\langle$P2]}}
                \texttt{SAKYTI: (1,5) yra tuščias langelis ('◌'), rėmelis yra ties (1,6). Patikrink dar kartą. Jei tau (1,5) netinka, siūlau I=(4,7) ir T=(6,7) arba I=(2,7) ir T=(1,7). Ką pasirinksi?} \\
            }
        }
    }
     \\ \\

    \theutterance \stepcounter{utterance}  
    & & & \multicolumn{2}{p{0.3\linewidth}}{
        \cellcolor[rgb]{0.9,0.9,0.9}{
            \makecell[{{p{\linewidth}}}]{
                \texttt{\tiny{[GM$|$GM]}}
                \texttt{SAKYTI: (1,5) yra tuščias langelis ('◌'), rėmelis yra ties (1,6). Patikrink dar kartą. Jei tau (1,5) netinka, siūlau I=(4,7) ir T=(6,7) arba I=(2,7) ir T=(1,7). Ką pasirinksi?} \\
            }
        }
    }
    & & \\ \\

    \theutterance \stepcounter{utterance}  
    & & \multicolumn{4}{p{0.6\linewidth}}{
        \cellcolor[rgb]{0.9,0.9,0.9}{
            \makecell[{{p{\linewidth}}}]{
                \texttt{\tiny{[P1$\langle$GM]}}
                \texttt{Jūsų žinutė buvo perduota kitam žaidėjui.} \\
\\ 
\\ 
\texttt{Šiuo metu žaidžiate 5 raundą iš maksimaliai 20.} \\
\texttt{Jūs bendrai sukaupėte 2 iš 12 baudų.} \\
\texttt{Kitas žaidėjas atsiuntė šią žinutę:} \\
\\ 
\\ 
\texttt{"(1,5) yra tuščias langelis ('◌'), rėmelis yra ties (1,6). Patikrink dar kartą. Jei tau (1,5) netinka, siūlau I=(4,7) ir T=(6,7) arba I=(2,7) ir T=(1,7). Ką pasirinksi?"} \\
\\ 
\\ 
\texttt{Kokia jūsų kita komanda?} \\
            }
        }
    }
    & \\ \\

    \theutterance \stepcounter{utterance}  
    & \multicolumn{4}{p{0.6\linewidth}}{
        \cellcolor[rgb]{0.9,0.9,0.9}{
            \makecell[{{p{\linewidth}}}]{
                \texttt{\tiny{[P1$\rangle$GM]}}
                \texttt{SAKYTI: Pasirenku I=(4,7) ir T=(6,7). Tinka tau?} \\
            }
        }
    }
    & & \\ \\

    \theutterance \stepcounter{utterance}  
    & & & \multicolumn{2}{p{0.3\linewidth}}{
        \cellcolor[rgb]{0.9,0.9,0.9}{
            \makecell[{{p{\linewidth}}}]{
                \texttt{\tiny{[GM$|$GM]}}
                \texttt{SAKYTI: Pasirenku I=(4,7) ir T=(6,7). Tinka tau?} \\
            }
        }
    }
    & & \\ \\

    \theutterance \stepcounter{utterance}  
    & & \multicolumn{4}{p{0.6\linewidth}}{
        \cellcolor[rgb]{0.9,0.9,0.9}{
            \makecell[{{p{\linewidth}}}]{
                \texttt{\tiny{[GM$\rangle$P2]}}
                \texttt{Jūsų žinutė buvo perduota kitam žaidėjui.} \\
\\ 
\\ 
\texttt{Šiuo metu žaidžiate 6 raundą iš maksimaliai 20.} \\
\texttt{Jūs bendrai sukaupėte 2 iš 12 baudų.} \\
\texttt{Kitas žaidėjas atsiuntė šią žinutę:} \\
\\ 
\\ 
\texttt{"Pasirenku I=(4,7) ir T=(6,7). Tinka tau?"} \\
\\ 
\\ 
\texttt{Kokia jūsų kita komanda?} \\
            }
        }
    }
    & \\ \\

    \theutterance \stepcounter{utterance}  
    & & & \multicolumn{4}{p{0.6\linewidth}}{
        \cellcolor[rgb]{0.9,0.9,0.9}{
            \makecell[{{p{\linewidth}}}]{
                \texttt{\tiny{[GM$\langle$P2]}}
                \texttt{SAKYTI: Tinka! Susitariame: H=(1,2), I=(4,7), T=(6,7), C=(3,4), W=(5,5). Dabar kelsiu I į (4,7).} \\
            }
        }
    }
     \\ \\

    \theutterance \stepcounter{utterance}  
    & & & \multicolumn{2}{p{0.3\linewidth}}{
        \cellcolor[rgb]{0.9,0.9,0.9}{
            \makecell[{{p{\linewidth}}}]{
                \texttt{\tiny{[GM$|$GM]}}
                \texttt{SAKYTI: Tinka! Susitariame: H=(1,2), I=(4,7), T=(6,7), C=(3,4), W=(5,5). Dabar kelsiu I į (4,7).} \\
            }
        }
    }
    & & \\ \\

    \theutterance \stepcounter{utterance}  
    & & \multicolumn{4}{p{0.6\linewidth}}{
        \cellcolor[rgb]{0.9,0.9,0.9}{
            \makecell[{{p{\linewidth}}}]{
                \texttt{\tiny{[P1$\langle$GM]}}
                \texttt{Jūsų žinutė buvo perduota kitam žaidėjui.} \\
\\ 
\\ 
\texttt{Šiuo metu žaidžiate 6 raundą iš maksimaliai 20.} \\
\texttt{Jūs bendrai sukaupėte 2 iš 12 baudų.} \\
\texttt{Kitas žaidėjas atsiuntė šią žinutę:} \\
\\ 
\\ 
\texttt{"Tinka! Susitariame: H=(1,2), I=(4,7), T=(6,7), C=(3,4), W=(5,5). Dabar kelsiu I į (4,7)."} \\
\\ 
\\ 
\texttt{Kokia jūsų kita komanda?} \\
            }
        }
    }
    & \\ \\

    \theutterance \stepcounter{utterance}  
    & \multicolumn{4}{p{0.6\linewidth}}{
        \cellcolor[rgb]{0.9,0.9,0.9}{
            \makecell[{{p{\linewidth}}}]{
                \texttt{\tiny{[P1$\rangle$GM]}}
                \texttt{JUDĖTI: I, (4,7)} \\
            }
        }
    }
    & & \\ \\

    \theutterance \stepcounter{utterance}  
    & & & \multicolumn{2}{p{0.3\linewidth}}{
        \cellcolor[rgb]{0.9,0.9,0.9}{
            \makecell[{{p{\linewidth}}}]{
                \texttt{\tiny{[GM$|$GM]}}
                \texttt{JUDĖTI: I, (4,7)} \\
            }
        }
    }
    & & \\ \\

    \theutterance \stepcounter{utterance}  
    & & & \multicolumn{2}{p{0.3\linewidth}}{
        \cellcolor[rgb]{0.9,0.9,0.9}{
            \makecell[{{p{\linewidth}}}]{
                \texttt{\tiny{[GM$|$GM]}}
                \texttt{I moved to (4, 7)} \\
            }
        }
    }
    & & \\ \\

    \theutterance \stepcounter{utterance}  
    & & \multicolumn{4}{p{0.6\linewidth}}{
        \cellcolor[rgb]{0.9,0.9,0.9}{
            \makecell[{{p{\linewidth}}}]{
                \texttt{\tiny{[GM$\rangle$P2]}}
                \texttt{Jūsų žinutė buvo perduota kitam žaidėjui.} \\
\\ 
\\ 
\texttt{Šiuo metu žaidžiate 7 raundą iš maksimaliai 20.} \\
\texttt{Jūs bendrai sukaupėte 2 iš 12 baudų.} \\
\texttt{Kitas žaidėjas perkėlė objektą savo tinklelyje.} \\
\\ 
\\ 
\texttt{Kokia jūsų kita komanda?} \\
            }
        }
    }
    & \\ \\

    \theutterance \stepcounter{utterance}  
    & & & \multicolumn{4}{p{0.6\linewidth}}{
        \cellcolor[rgb]{0.9,0.9,0.9}{
            \makecell[{{p{\linewidth}}}]{
                \texttt{\tiny{[GM$\langle$P2]}}
                \texttt{JUDĖTI: I, (4, 7)} \\
            }
        }
    }
     \\ \\

    \theutterance \stepcounter{utterance}  
    & & & \multicolumn{2}{p{0.3\linewidth}}{
        \cellcolor[rgb]{0.9,0.9,0.9}{
            \makecell[{{p{\linewidth}}}]{
                \texttt{\tiny{[GM$|$GM]}}
                \texttt{JUDĖTI: I, (4, 7)} \\
            }
        }
    }
    & & \\ \\

    \theutterance \stepcounter{utterance}  
    & & & \multicolumn{2}{p{0.3\linewidth}}{
        \cellcolor[rgb]{0.9,0.9,0.9}{
            \makecell[{{p{\linewidth}}}]{
                \texttt{\tiny{[GM$|$GM]}}
                \texttt{Nuobauda: (4,7) nėra tuščias, bet jame yra 'I'.} \\
\\ 
\texttt{Jūs bendrai sukaupėte 3 iš 12 baudų.} \\
\texttt{Pabandykite dar kartą!} \\
            }
        }
    }
    & & \\ \\

    \theutterance \stepcounter{utterance}  
    & & \multicolumn{4}{p{0.6\linewidth}}{
        \cellcolor[rgb]{0.9,0.9,0.9}{
            \makecell[{{p{\linewidth}}}]{
                \texttt{\tiny{[GM$\rangle$P2]}}
                \texttt{Nuobauda: (4,7) nėra tuščias, bet jame yra 'I'.} \\
\\ 
\texttt{Jūs bendrai sukaupėte 3 iš 12 baudų.} \\
\texttt{Pabandykite dar kartą!} \\
            }
        }
    }
    & \\ \\

    \theutterance \stepcounter{utterance}  
    & & & \multicolumn{4}{p{0.6\linewidth}}{
        \cellcolor[rgb]{0.9,0.9,0.9}{
            \makecell[{{p{\linewidth}}}]{
                \texttt{\tiny{[GM$\langle$P2]}}
                \texttt{Ačiū, teisingai {-} I jau yra (4,7). Tada kelsiu T į (6,7).} \\
\\ 
\texttt{JUDĖTI: T, (6, 7)} \\
            }
        }
    }
     \\ \\

    \theutterance \stepcounter{utterance}  
    & & & \multicolumn{2}{p{0.3\linewidth}}{
        \cellcolor[rgb]{0.9,0.9,0.9}{
            \makecell[{{p{\linewidth}}}]{
                \texttt{\tiny{[GM$|$GM]}}
                \texttt{Ačiū, teisingai {-} I jau yra (4,7). Tada kelsiu T į (6,7).} \\
\\ 
\texttt{JUDĖTI: T, (6, 7)} \\
            }
        }
    }
    & & \\ \\

    \theutterance \stepcounter{utterance}  
    & & & \multicolumn{2}{p{0.3\linewidth}}{
        \cellcolor[rgb]{0.9,0.9,0.9}{
            \makecell[{{p{\linewidth}}}]{
                \texttt{\tiny{[GM$|$GM]}}
                \texttt{Invalid format: head is not empty: 'Ačiū, teisingai {-} I jau yra (4,7). Tada kelsiu T į (6,7).} \\
\\ 
\texttt{'} \\
            }
        }
    }
    & & \\ \\

    \theutterance \stepcounter{utterance}  
    & & \multicolumn{4}{p{0.6\linewidth}}{
        \cellcolor[rgb]{0.9,0.9,0.9}{
            \makecell[{{p{\linewidth}}}]{
                \texttt{\tiny{[GM$\rangle$P2]}}
                \texttt{Nuobauda: Jūsų žinutėje neturi būti nieko prieš komandą!} \\
\\ 
\texttt{Įsitikinkite, kad jūsų atsakyme yra tik `SAKYTI: <ŽINUTĖ>` arba `JUDĖTI: <OBJEKTAS>, (<X>, <Y>)`, ir nieko daugiau!} \\
\\ 
\texttt{Jūs bendrai sukaupėte 4 iš 12 baudų.} \\
\texttt{Pabandykite dar kartą!} \\
            }
        }
    }
    & \\ \\

    \theutterance \stepcounter{utterance}  
    & & & \multicolumn{4}{p{0.6\linewidth}}{
        \cellcolor[rgb]{0.9,0.9,0.9}{
            \makecell[{{p{\linewidth}}}]{
                \texttt{\tiny{[GM$\langle$P2]}}
                \texttt{JUDĖTI: T, (6, 7)} \\
            }
        }
    }
     \\ \\

    \theutterance \stepcounter{utterance}  
    & & & \multicolumn{2}{p{0.3\linewidth}}{
        \cellcolor[rgb]{0.9,0.9,0.9}{
            \makecell[{{p{\linewidth}}}]{
                \texttt{\tiny{[GM$|$GM]}}
                \texttt{JUDĖTI: T, (6, 7)} \\
            }
        }
    }
    & & \\ \\

    \theutterance \stepcounter{utterance}  
    & & & \multicolumn{2}{p{0.3\linewidth}}{
        \cellcolor[rgb]{0.9,0.9,0.9}{
            \makecell[{{p{\linewidth}}}]{
                \texttt{\tiny{[GM$|$GM]}}
                \texttt{T moved to (6, 7)} \\
            }
        }
    }
    & & \\ \\

    \theutterance \stepcounter{utterance}  
    & & \multicolumn{4}{p{0.6\linewidth}}{
        \cellcolor[rgb]{0.9,0.9,0.9}{
            \makecell[{{p{\linewidth}}}]{
                \texttt{\tiny{[P1$\langle$GM]}}
                \texttt{Sėkmingai perkelta 'I' į (4,7). Jūsų atnaujintas tinklelis atrodo taip:} \\
\texttt{```} \\
\texttt{ 1234567} \\
\texttt{╔═╤═╤═══╗} \\
\texttt{║C│◌│◌T◌║ 1} \\
\texttt{║H├─┤◌◌◌║ 2} \\
\texttt{╟─┤◌├─┬─╢ 3} \\
\texttt{║◌│◌│◌│◌║ 4} \\
\texttt{╟─┼─┘W│◌║ 5} \\
\texttt{║◌│◌◌┌┴─╢ 6} \\
\texttt{║◌│◌I│◌◌║ 7} \\
\texttt{╚═╧══╧══╝} \\
\texttt{```} \\
\texttt{Šiuo metu žaidžiate 7 raundą iš maksimaliai 20.} \\
\texttt{Jūs bendrai sukaupėte 4 iš 12 baudų.} \\
\texttt{Kitas žaidėjas perkėlė objektą savo tinklelyje.} \\
\\ 
\\ 
\texttt{Kokia jūsų kita komanda?} \\
            }
        }
    }
    & \\ \\

    \theutterance \stepcounter{utterance}  
    & \multicolumn{4}{p{0.6\linewidth}}{
        \cellcolor[rgb]{0.9,0.9,0.9}{
            \makecell[{{p{\linewidth}}}]{
                \texttt{\tiny{[P1$\rangle$GM]}}
                \texttt{JUDĖTI: T, (6,7)} \\
            }
        }
    }
    & & \\ \\

    \theutterance \stepcounter{utterance}  
    & & & \multicolumn{2}{p{0.3\linewidth}}{
        \cellcolor[rgb]{0.9,0.9,0.9}{
            \makecell[{{p{\linewidth}}}]{
                \texttt{\tiny{[GM$|$GM]}}
                \texttt{JUDĖTI: T, (6,7)} \\
            }
        }
    }
    & & \\ \\

    \theutterance \stepcounter{utterance}  
    & & & \multicolumn{2}{p{0.3\linewidth}}{
        \cellcolor[rgb]{0.9,0.9,0.9}{
            \makecell[{{p{\linewidth}}}]{
                \texttt{\tiny{[GM$|$GM]}}
                \texttt{T moved to (6, 7)} \\
            }
        }
    }
    & & \\ \\

    \theutterance \stepcounter{utterance}  
    & & \multicolumn{4}{p{0.6\linewidth}}{
        \cellcolor[rgb]{0.9,0.9,0.9}{
            \makecell[{{p{\linewidth}}}]{
                \texttt{\tiny{[GM$\rangle$P2]}}
                \texttt{Sėkmingai perkelta 'T' į (6,7). Jūsų atnaujintas tinklelis atrodo taip:} \\
\texttt{```} \\
\texttt{ 1234567} \\
\texttt{╔═╤═╤═══╗} \\
\texttt{║◌│◌│◌◌◌║ 1} \\
\texttt{║H├─┤◌◌◌║ 2} \\
\texttt{╟─┤◌├─┬─╢ 3} \\
\texttt{║◌│◌│C│◌║ 4} \\
\texttt{╟─┼─┘W│◌║ 5} \\
\texttt{║◌│◌◌┌┴─╢ 6} \\
\texttt{║◌│◌I│T◌║ 7} \\
\texttt{╚═╧══╧══╝} \\
\texttt{```} \\
\texttt{Šiuo metu žaidžiate 8 raundą iš maksimaliai 20.} \\
\texttt{Jūs bendrai sukaupėte 4 iš 12 baudų.} \\
\texttt{Kitas žaidėjas perkėlė objektą savo tinklelyje.} \\
\\ 
\\ 
\texttt{Kokia jūsų kita komanda?} \\
            }
        }
    }
    & \\ \\

    \theutterance \stepcounter{utterance}  
    & & & \multicolumn{4}{p{0.6\linewidth}}{
        \cellcolor[rgb]{0.9,0.9,0.9}{
            \makecell[{{p{\linewidth}}}]{
                \texttt{\tiny{[GM$\langle$P2]}}
                \texttt{SAKYTI: Aš perkėliau T į (6,7). Mano lentoje dabar visi objektai tikslinėse pozicijose: H=(1,2), I=(4,7), T=(6,7), C=(3,4), W=(5,5). Ar tu pasiruošęs baigti?} \\
            }
        }
    }
     \\ \\

    \theutterance \stepcounter{utterance}  
    & & & \multicolumn{2}{p{0.3\linewidth}}{
        \cellcolor[rgb]{0.9,0.9,0.9}{
            \makecell[{{p{\linewidth}}}]{
                \texttt{\tiny{[GM$|$GM]}}
                \texttt{SAKYTI: Aš perkėliau T į (6,7). Mano lentoje dabar visi objektai tikslinėse pozicijose: H=(1,2), I=(4,7), T=(6,7), C=(3,4), W=(5,5). Ar tu pasiruošęs baigti?} \\
            }
        }
    }
    & & \\ \\

    \theutterance \stepcounter{utterance}  
    & & \multicolumn{4}{p{0.6\linewidth}}{
        \cellcolor[rgb]{0.9,0.9,0.9}{
            \makecell[{{p{\linewidth}}}]{
                \texttt{\tiny{[P1$\langle$GM]}}
                \texttt{Sėkmingai perkelta 'T' į (6,7). Jūsų atnaujintas tinklelis atrodo taip:} \\
\texttt{```} \\
\texttt{ 1234567} \\
\texttt{╔═╤═╤═══╗} \\
\texttt{║C│◌│◌◌◌║ 1} \\
\texttt{║H├─┤◌◌◌║ 2} \\
\texttt{╟─┤◌├─┬─╢ 3} \\
\texttt{║◌│◌│◌│◌║ 4} \\
\texttt{╟─┼─┘W│◌║ 5} \\
\texttt{║◌│◌◌┌┴─╢ 6} \\
\texttt{║◌│◌I│T◌║ 7} \\
\texttt{╚═╧══╧══╝} \\
\texttt{```} \\
\texttt{Šiuo metu žaidžiate 8 raundą iš maksimaliai 20.} \\
\texttt{Jūs bendrai sukaupėte 4 iš 12 baudų.} \\
\texttt{Kitas žaidėjas atsiuntė šią žinutę:} \\
\\ 
\\ 
\texttt{"Aš perkėliau T į (6,7). Mano lentoje dabar visi objektai tikslinėse pozicijose: H=(1,2), I=(4,7), T=(6,7), C=(3,4), W=(5,5). Ar tu pasiruošęs baigti?"} \\
\\ 
\\ 
\texttt{Kokia jūsų kita komanda?} \\
            }
        }
    }
    & \\ \\

    \theutterance \stepcounter{utterance}  
    & \multicolumn{4}{p{0.6\linewidth}}{
        \cellcolor[rgb]{0.9,0.9,0.9}{
            \makecell[{{p{\linewidth}}}]{
                \texttt{\tiny{[P1$\rangle$GM]}}
                \texttt{JUDĖTI: C, (3,4)} \\
            }
        }
    }
    & & \\ \\

    \theutterance \stepcounter{utterance}  
    & & & \multicolumn{2}{p{0.3\linewidth}}{
        \cellcolor[rgb]{0.9,0.9,0.9}{
            \makecell[{{p{\linewidth}}}]{
                \texttt{\tiny{[GM$|$GM]}}
                \texttt{JUDĖTI: C, (3,4)} \\
            }
        }
    }
    & & \\ \\

    \theutterance \stepcounter{utterance}  
    & & & \multicolumn{2}{p{0.3\linewidth}}{
        \cellcolor[rgb]{0.9,0.9,0.9}{
            \makecell[{{p{\linewidth}}}]{
                \texttt{\tiny{[GM$|$GM]}}
                \texttt{C moved to (3, 4)} \\
            }
        }
    }
    & & \\ \\

    \theutterance \stepcounter{utterance}  
    & & \multicolumn{4}{p{0.6\linewidth}}{
        \cellcolor[rgb]{0.9,0.9,0.9}{
            \makecell[{{p{\linewidth}}}]{
                \texttt{\tiny{[GM$\rangle$P2]}}
                \texttt{Jūsų žinutė buvo perduota kitam žaidėjui.} \\
\\ 
\\ 
\texttt{Šiuo metu žaidžiate 9 raundą iš maksimaliai 20.} \\
\texttt{Jūs bendrai sukaupėte 4 iš 12 baudų.} \\
\texttt{Kitas žaidėjas perkėlė objektą savo tinklelyje.} \\
\\ 
\\ 
\texttt{Kokia jūsų kita komanda?} \\
            }
        }
    }
    & \\ \\

    \theutterance \stepcounter{utterance}  
    & & & \multicolumn{4}{p{0.6\linewidth}}{
        \cellcolor[rgb]{0.9,0.9,0.9}{
            \makecell[{{p{\linewidth}}}]{
                \texttt{\tiny{[GM$\langle$P2]}}
                \texttt{SAKYTI: baigta?} \\
            }
        }
    }
     \\ \\

    \theutterance \stepcounter{utterance}  
    & & & \multicolumn{2}{p{0.3\linewidth}}{
        \cellcolor[rgb]{0.9,0.9,0.9}{
            \makecell[{{p{\linewidth}}}]{
                \texttt{\tiny{[GM$|$GM]}}
                \texttt{SAKYTI: baigta?} \\
            }
        }
    }
    & & \\ \\

    \theutterance \stepcounter{utterance}  
    & & \multicolumn{4}{p{0.6\linewidth}}{
        \cellcolor[rgb]{0.9,0.9,0.9}{
            \makecell[{{p{\linewidth}}}]{
                \texttt{\tiny{[P1$\langle$GM]}}
                \texttt{Sėkmingai perkelta 'C' į (3,4). Jūsų atnaujintas tinklelis atrodo taip:} \\
\texttt{```} \\
\texttt{ 1234567} \\
\texttt{╔═╤═╤═══╗} \\
\texttt{║◌│◌│◌◌◌║ 1} \\
\texttt{║H├─┤◌◌◌║ 2} \\
\texttt{╟─┤◌├─┬─╢ 3} \\
\texttt{║◌│C│◌│◌║ 4} \\
\texttt{╟─┼─┘W│◌║ 5} \\
\texttt{║◌│◌◌┌┴─╢ 6} \\
\texttt{║◌│◌I│T◌║ 7} \\
\texttt{╚═╧══╧══╝} \\
\texttt{```} \\
\texttt{Šiuo metu žaidžiate 9 raundą iš maksimaliai 20.} \\
\texttt{Jūs bendrai sukaupėte 4 iš 12 baudų.} \\
\texttt{Kitas žaidėjas atsiuntė šią žinutę:} \\
\\ 
\\ 
\texttt{"baigta?"} \\
\\ 
\\ 
\texttt{Kokia jūsų kita komanda?} \\
            }
        }
    }
    & \\ \\

    \theutterance \stepcounter{utterance}  
    & \multicolumn{4}{p{0.6\linewidth}}{
        \cellcolor[rgb]{0.9,0.9,0.9}{
            \makecell[{{p{\linewidth}}}]{
                \texttt{\tiny{[P1$\rangle$GM]}}
                \texttt{SAKYTI: baigta!} \\
            }
        }
    }
    & & \\ \\

    \theutterance \stepcounter{utterance}  
    & & & \multicolumn{2}{p{0.3\linewidth}}{
        \cellcolor[rgb]{0.9,0.9,0.9}{
            \makecell[{{p{\linewidth}}}]{
                \texttt{\tiny{[GM$|$GM]}}
                \texttt{SAKYTI: baigta!} \\
            }
        }
    }
    & & \\ \\

    \theutterance \stepcounter{utterance}  
    & & & \multicolumn{2}{p{0.3\linewidth}}{
        \cellcolor[rgb]{0.9,0.9,0.9}{
            \makecell[{{p{\linewidth}}}]{
                \texttt{\tiny{[GM$|$GM]}}
                \texttt{true} \\
            }
        }
    }
    & & \\ \\

    \theutterance \stepcounter{utterance}  
    & & & \multicolumn{2}{p{0.3\linewidth}}{
        \cellcolor[rgb]{0.9,0.9,0.9}{
            \makecell[{{p{\linewidth}}}]{
                \texttt{\tiny{[GM$|$GM]}}
                \texttt{Initial states:} \\
\texttt{```} \\
\texttt{Player 1:} \\
\texttt{ 1234567} \\
\texttt{╔═╤═╤═══╗} \\
\texttt{║C│◌│◌TI║ 1} \\
\texttt{║◌├─┤W◌◌║ 2} \\
\texttt{╟─┤◌├─┬─╢ 3} \\
\texttt{║◌│◌│◌│◌║ 4} \\
\texttt{╟─┼─┘◌│◌║ 5} \\
\texttt{║◌│◌H┌┴─╢ 6} \\
\texttt{║◌│◌◌│◌◌║ 7} \\
\texttt{╚═╧══╧══╝} \\
\\ 
\\ 
\texttt{Player 2:} \\
\texttt{ 1234567} \\
\texttt{╔═╤═╤═══╗} \\
\texttt{║H│◌│◌◌◌║ 1} \\
\texttt{║◌├─┤◌◌◌║ 2} \\
\texttt{╟─┤◌├─┬─╢ 3} \\
\texttt{║◌│◌│C│◌║ 4} \\
\texttt{╟─┼─┘◌│◌║ 5} \\
\texttt{║◌│◌◌┌┴─╢ 6} \\
\texttt{║T│WI│◌◌║ 7} \\
\texttt{╚═╧══╧══╝} \\
\\ 
\texttt{```} \\
            }
        }
    }
    & & \\ \\

    \theutterance \stepcounter{utterance}  
    & & & \multicolumn{2}{p{0.3\linewidth}}{
        \cellcolor[rgb]{0.9,0.9,0.9}{
            \makecell[{{p{\linewidth}}}]{
                \texttt{\tiny{[GM$|$GM]}}
                \texttt{End states:} \\
\texttt{```} \\
\texttt{Player 1:} \\
\texttt{ 1234567} \\
\texttt{╔═╤═╤═══╗} \\
\texttt{║◌│◌│◌◌◌║ 1} \\
\texttt{║H├─┤◌◌◌║ 2} \\
\texttt{╟─┤◌├─┬─╢ 3} \\
\texttt{║◌│C│◌│◌║ 4} \\
\texttt{╟─┼─┘W│◌║ 5} \\
\texttt{║◌│◌◌┌┴─╢ 6} \\
\texttt{║◌│◌I│T◌║ 7} \\
\texttt{╚═╧══╧══╝} \\
\\ 
\\ 
\texttt{Player 2:} \\
\texttt{ 1234567} \\
\texttt{╔═╤═╤═══╗} \\
\texttt{║◌│◌│◌◌◌║ 1} \\
\texttt{║H├─┤◌◌◌║ 2} \\
\texttt{╟─┤◌├─┬─╢ 3} \\
\texttt{║◌│◌│C│◌║ 4} \\
\texttt{╟─┼─┘W│◌║ 5} \\
\texttt{║◌│◌◌┌┴─╢ 6} \\
\texttt{║◌│◌I│T◌║ 7} \\
\texttt{╚═╧══╧══╝} \\
\\ 
\texttt{```} \\
            }
        }
    }
    & & \\ \\

    \theutterance \stepcounter{utterance}  
    & & & \multicolumn{2}{p{0.3\linewidth}}{
        \cellcolor[rgb]{0.9,0.9,0.9}{
            \makecell[{{p{\linewidth}}}]{
                \texttt{\tiny{[GM$|$GM]}}
                \texttt{* success: True} \\
\texttt{* lose: False} \\
\texttt{* aborted: False} \\
\texttt{{-}{-}{-}{-}{-}{-}{-}} \\
\texttt{* Shifts: 4.00} \\
\texttt{* Max Shifts: 8.00} \\
\texttt{* Min Shifts: 4.00} \\
\texttt{* End Distance Sum: 2.00} \\
\texttt{* Init Distance Sum: 30.73} \\
\texttt{* Expected Distance Sum: 20.95} \\
\texttt{* Penalties: 4.00} \\
\texttt{* Max Penalties: 12.00} \\
\texttt{* Rounds: 9.00} \\
\texttt{* Max Rounds: 20.00} \\
\texttt{* Object Count: 5.00} \\
            }
        }
    }
    & & \\ \\

\end{supertabular}
}

\end{document}
